%%%%%%%%%%%%%%%%%%%%%%%%%%%%%%%%%%%%%%%%%%%%%%%%%%%%%%%
% % Lines starting with % are comments, which are ignored.
% % This is a handy way of indicating the date and version of
% % your document, to wit:
% %
% % LaTeX sample file
% % Modified March, 2002
% %
%%%%%%%%%%%%%%%%%%%%%%%%%%%%%%%%%%%%%%%%%%%%%%%%%%%%%%%
\DocumentMetadata{
  % PDF standards: ONE key whose value is a comma-separated list --
  % UA-2 (accessibility) plus A-4f (archival; the archival companion of
  % UA-2). Do NOT use separate pdfstandard=... keys: with ordinary
  % key-value parsing only the last one takes effect.
  pdfstandard={UA-2,A-4f},
  lang=en-US,
  tagging=on,
  tagging-setup={
    math/setup=mathml-SE
  }
}

\documentclass[12pt]{exam}
% worksheet-macros.tex -- shared package setup for worksheet-template.tex.
%
% NOTE ON PATHS: every template in this repository is stored one level up
% from where it is compiled (see README.md). All relative \input paths
% below are resolved against the folder the document is compiled in,
% which sits one directory above this file.

% --- Math ------------------------------------------------------------
% amsmath: equation environments and numbering
% amsfonts: \mathbb etc.  |  mathrsfs: \mathfrak  |  stmaryrd: extra symbols
\usepackage{amsmath}
\usepackage{amsfonts}
\usepackage{mathrsfs}
\usepackage{stmaryrd}

% --- Layout and glyphs ------------------------------------------------
% float: [H] float placement (stop floats where you put them)
% fontspec: Unicode/OpenType fonts (compile with XeLaTeX or LuaLaTeX)
% emoji: emoji glyphs in the text
\usepackage{float}
\usepackage{fontspec}
\usepackage{emoji}

% --- Typography --------------------------------------------------------
% Atkinson Hyperlegible as the document-wide sans serif (sfdefault):
% a high-legibility face that is easier to read for students with
% dyslexia or low vision.
\usepackage[sfdefault]{atkinson}

% --- Colors -------------------------------------------------------------
% dvipsnames: extended named color palette (the Accessible* colors in
%   lecture-macros.tex are built on it);
% table: \rowcolors / \cellcolor for shading worksheet tables.
% NOTE: xcolor may only be loaded ONCE per document -- a second
% \usepackage{xcolor} with different options is a fatal "Option clash"
% error, so ALL of its options are listed here.
\usepackage[dvipsnames,table]{xcolor}

% --- Multimedia ---------------------------------------------------------
% animate: embeds audio/video (and legacy Flash) media in the worksheet.
\usepackage{animate}

% --- Bibliography --------------------------------------------------------
% biblatex (biber backend): citations and a reference list for worksheets.
\usepackage[
  backend=biber,
  style=alphabetic,
  url=true,
  doi=true,
  isbn=true,
  backref=true
]{biblatex}

% --- Links ---------------------------------------------------------------
% hyperref makes every \ref/citation/URL clickable in the PDF.
% Link colors must satisfy WCAG 1.4.3 (>= 4.5:1 contrast on white):
%   blue = 8.6:1, green!40!black = 7.2:1.
% (The old cyan/magenta values measured 1.25:1 / 3.1:1 and failed.)
\usepackage{hyperref}
\hypersetup{
  colorlinks=true,
  linkcolor=blue,
  filecolor=blue,
  urlcolor=green!40!black,
}

% --- Shared course macros ------------------------------------------------
% Root macros.tex: math shorthands used across all course documents
% (\RR, \NN, \ZZ, \QQ, \grad, \curl, \tr, \der, \pder, ...).
\newcommand{\RR}{\mathbb{R}}
\newcommand{\NN}{\mathbb{N}}
\newcommand{\ZZ}{\mathbb{Z}}
\newcommand{\QQ}{\mathbb{Q}}
\newcommand{\vv}{\vec{v}}
\newcommand{\ww}{\vec{w}}
\newcommand{\xx}{\vec{x}}
\newcommand{\yy}{\vec{y}}
\newcommand{\zz}{\vec{z}}
\newcommand{\uu}{\vec{u}}
\newcommand{\basis}[1]{\vec{e_{#1}}}
\newcommand{\id}{\mathrm{id}}
\newcommand{\CC}{\mathbb{C}}
\renewcommand{\div}{\operatorname{div}}
\newcommand{\grad}{\operatorname{grad}}
\newcommand{\curl}{\operatorname{curl}}
\newcommand{\tr}{\operatorname{tr}}
\newcommand{\der}[2]{\dfrac{d #1}{d #2}}
\newcommand{\pder}[3][]{\dfrac{\partial^{#1} #2}{\partial #3}}


% % Title and author(s)
%%%%%%%%%%%%%%%%%%%%%%%%%%%%%%%%%%%%%%%%%%%%%%%%%%%%%%%
\title{Worksheet 11 - The Chain Rule}
\author{MATH 1271}
\date{10/15/2026}
%%%%%%%%%%%%%%%%%%%%%%%%%%%%%%%%%%%%%%%%%%%%%%%%%%%%%%%
% %
% % The next command allows your in import encapsulated
% % postscript files, .epsf or .eps files, which
% % contain vector graphic image data.
% %
%%%%%%%%%%%%%%%%%%%%%%%%%%%%%%%%%%%%%%%%%%%%%%%%%%%%%%%
% %   The first thanks indicates your affiliation
% %
% %  Just the name here.
% %
% % Your mailing address goes at the end.
% %
% % \thanks is also how you indicate grant support
% %
%%%%%%%%%%%%%%%%%%%%%%%%%%%%%%%%%%%%%%%%%%%%%%%%%%%%%%%

%%%%%%%%%%%%%%%%%%%%%%%%%%%%%%%%%%%%%%%%%%%%%%%%%%%%%%%
% $ The next command ensures the margins are 1 inch on
% % all sides.
% %
%%%%%%%%%%%%%%%%%%%%%%%%%%%%%%%%%%%%%%%%%%%%%%%%%%%%%%%
\begin{document}
\newpage
\maketitle

\begin{enumerate}

\item For each function below, identify how to express it as a composite of elementary functions and then find its derivative:

\begin{enumerate}


\item $m(v)=\sin(v^2)\cos(v^3)$

\vfill

\item $h(y)=\dfrac{\cos(10y)}{e^{4y}+1}$

\vfill

\item $g(w) = (6w^2 + 7w)^4$

  \vfill

\item $\alpha(x) = \left(\sqrt[3]{12e^{4x} + 3e^{4x} + 2} + \sin^2(e^{4x})\right)^{-1}$
  \vfill
\end{enumerate}

\newpage


\item Use known derivative rules, including the Chain Rule, as needed to answer each of the following questions:

\begin{enumerate}

\item Find an equation for the tangent line to the curve $y=\sqrt{e^x+3}$ at the point where $x=0$.

\vfill

\item If $s(t)=\frac{1}{(t^2+1)^3}$ represents the position function of a particle moving horizontally along an axis at time $t$, find the particle's instantaneous velocity at $t=1$. Is the particle moving to the left or to the right at that instance? (Hint: What is the sign of $v(1)$?)

\vfill

\item At sea level, the air pressure is $30$ inches of mercury. At an altitude of $h$ feet above sea level, the air pressure, $P$, is given by the function $P=30e^{-0.0000323h}$. Compute $\frac{dP}{dh}$ and explain what this derivative function tells you about air pressure, including a discussion of units of $\frac{dP}{dh}$.

\vfill


\newpage

\item Referring to the previous question, determine how fast the air pressure is changing for a pilot of a small plane passing through an altitude of $1,000$ feet. (Round your answer to $6$ decimal places.)

\vfill

\item Suppose that $f(x)$ and $g(x)$ are differentiable functions and that the following information about them is known:

\begin{center}
\begin{tabular}{c | c | c | c | c}
$x$ & $f(x)$ & $f'(x)$ & $g(x)$ & $g'(x)$\\
\hline
$-1$ & $2$ & $-5$ & $-3$ & $4$\\
\hline
$2$ & $-3$ & $4$ & $-1$ & $2$
\end{tabular}
\end{center}

If $C(x)$ is a function given by the formula $C(x)=f(g(x))$, determine $C'(2)$.

\vfill

\item Using the same table as in the previous question, let $D(x)$ be the function $f(f(x))$. Find $D'(-1)$.

\vfill




\end{enumerate}
\newpage

\item Say you are working for a company who produces solar panels, and are trying to determine
  the optimal density of photovoltaic cells with which to produce your panels.
  Your economists have determined that clients are willing to pay more for panels as
  the energy density (in watts per square meter) of your panels vary; 
  your chemists have determined that the energy density changes depending on the density of your
  photovoltaic cells, and that the cost per unit area product (dollars per square meter) 
  also only depends on that photovoltaic density (measured in cells per square meter).

  Suppose you know that:
    \begin{itemize}
      \item Your company currently produces panels with a PVD of 100 cells per square meter
      \item At a PVD of 100 cells per square meter, the energy density produced 175 watts per meter squared
      \item At that same PVD, panels cost about 60 dollars per square meter
      \item At an energy density of 175 watts per meter squared, clients are willing to pay 90 dollars per square meter
      \item At your current PVD, small changes to the PVD cause the energy density to change by about 0.02 watts per cell
      \item At your current PVD, small changes to it will increase the cost by about 0.20 USD per cell
      \item From market research, increasing energy density will make clients increase demand at a rate of about 2 dollars per watt.
    \end{itemize}

    \begin{enumerate}
      \item If the energy density as a function of photovoltaic density is given by $E(p)$, the cost as a function of photovoltaic density is given by $C(p)$, and the price a client is willing to pay as a function of energy density is given by $W(e)$, determine an expression
      for the profit per square meter after selling a solar panel.
      \vfill

      \item At our current values, how much profit are we making per square meter?
      \vfill
      \newpage

    \item Using the previously defined functions, rewrite the measurements on the previous page in the form of values of function/their derivatives.
      \vfill

    \item What is an expression for the derivative of profit as a function of photovoltaic density?
      \vfill

    \item Should you increase or decrease photovoltaic density if you want to increase profit?
      \vfill
    \end{enumerate}
\end{enumerate}
% Most people don't use these, so they are "commented out"
% by starting the lines with a "%"
%\begin{keywords}
%   \LaTeX, typesetting
%\end{keywords}

%\begin{AMS}
%   50C60, 18C25
%\end{AMS}

%%%%%%%%%%%%%%%%%%%%%%
% % Here is the start of the Text
%%%%%%%%%%%%%%%%%%%%%%
\end{document}
